\section{From the connection zero to resolvents and effective Hamiltonians}
\label{sec:scattering}
% BEGIN CURATED SUBJECT INDEX
\index{resolvent}
\index{effective Hamiltonian}
% END CURATED SUBJECT INDEX

\subsection{Lippmann--Schwinger equation and the transition operator}
\label{subsec:lippmann-schwinger}
% BEGIN CURATED SUBJECT INDEX
\index{operator!adjoint and self-adjoint}
\index{resolvent}
\index{Jost function}
\index{Lippmann--Schwinger equation}
\index{boundary condition!incoming and outgoing}
% END CURATED SUBJECT INDEX

Write the Hamiltonian as $H=H_0+V$, where $H_0$ is the asymptotic Hamiltonian
and $V$ is the interaction.  The resolvent identity gives
\begin{align}
  \Resolvent(z)
  &=\Resolvent_0(z)
  +\Resolvent_0(z)V\Resolvent(z),
  \label{eq:resolvent-identity}\\
  \Resolvent_0(z)&=(z-H_0)^{-1} .
\end{align}
For a generalized free eigenstate $\ket{\phi_E}$, the incoming and outgoing
scattering states are the boundary values
\begin{equation}
  \ket{\psi_E^{(\pm)}}
  =\ket{\phi_E}
  +\Resolvent_0(E\pm\rmi0)V\ket{\psi_E^{(\pm)}} .
  \label{eq:lippmann-schwinger}
\end{equation}
The transition operator $T(z)$ is defined by
\begin{equation}
  T(z)=V+V\Resolvent_0(z)T(z) .
  \label{eq:T-operator}
\end{equation}
On shell, matrix elements of $T(E+\rmi0)$ determine scattering amplitudes.  A
resonance appears as a pole of $T(z)$ and of the full resolvent.  The
Lippmann--Schwinger equation therefore expresses in operator language what the
Jost-function construction expressed through differential equations.

For real $E$ above threshold, the discontinuity of the free resolvent is
\begin{equation}
  \Resolvent_0(E+\rmi0)-\Resolvent_0(E-\rmi0)
  =-2\pi\rmi\,\delta(E-H_0) .
  \label{eq:resolvent-discontinuity}
\end{equation}
Together with Eq.~\eqref{eq:T-operator}, this identity yields the optical
theorem.  It also shows how an imaginary part can arise from an underlying
self-adjoint theory: the imaginary part is the boundary value of a continuum
resolvent, not a fundamental violation of probability conservation.

\subsection{Feshbach projection and energy-dependent effective Hamiltonian}
\label{subsec:feshbach}
% BEGIN CURATED SUBJECT INDEX
\index{operator!adjoint and self-adjoint}
\index{Feshbach projection}
\index{effective Hamiltonian}
\index{self-energy}
\index{exceptional point}
% END CURATED SUBJECT INDEX

Let $P$ be an orthogonal projection onto a selected subspace and let
$Q=\identity-P$.  Decompose a stationary state as
$\ket{\Psi}=P\ket{\Psi}+Q\ket{\Psi}$.  Eliminating the $Q$ component gives
\begin{equation}
  \left[z-\vsub{H}{eff}(z)\right]P\ket{\Psi}=0,
  \label{eq:feshbach-equation}
\end{equation}
with
\begin{equation}
  \vsub{H}{eff}(z)
  =PHP+PHQ\frac{1}{z-QHQ}QHP .
  \label{eq:feshbach-H}
\end{equation}
The second term is the self-energy.  For $z=E+\rmi0$ above a $Q$-space
threshold, it has the decomposition
\begin{equation}
  \Sigma(E+\rmi0)
  =\Delta(E)-\frac{\rmi}{2}W(E),
  \qquad W(E)\succeq0 .
  \label{eq:self-energy-decomposition}
\end{equation}
Here $\Delta(E)$ is a dispersive energy shift, $W(E)$ is a positive
semidefinite decay matrix, and $\succeq0$ denotes positive semidefiniteness.
Thus $\vsub{H}{eff}(E+\rmi0)$ is non-self-adjoint even though $H$ is
self-adjoint.  Its non-self-adjoint part records probability flux into the
eliminated continuum
\cite{Feshbach1958UnifiedI,Feshbach1962UnifiedII,Rotter:2007zng}.

The resonance equation is nonlinear because the operator depends on its own
spectral parameter:
\begin{equation}
  \det\left[z-\vsub{H}{eff}(z)\right]=0 .
  \label{eq:feshbach-pole-equation}
\end{equation}
Replacing $\vsub{H}{eff}(z)$ by a constant matrix near a narrow pole is a
controlled local approximation only if the self-energy varies slowly.  Near a
threshold, $\Sigma(z)$ has a branch point, and a finite constant matrix cannot
reproduce the correct nonanalyticity.  This observation is important for
exceptional points involving a continuum: the branch order can differ from
that of a generic finite-dimensional exceptional point
\cite{Garmon2021AnomalousEP}.

\subsection{Equivalence of the Feshbach and Siegert pictures}
\label{subsec:feshbach-siegert}
% BEGIN CURATED SUBJECT INDEX
\index{operator!adjoint and self-adjoint}
\index{unbounded operator}
\index{Hilbert space}
\index{boundary condition!incoming and outgoing}
\index{Gamow--Siegert state}
\index{exact WKB}
\index{rigged Hilbert space}
\index{Feshbach projection}
\index{self-energy}
% END CURATED SUBJECT INDEX

The Feshbach and outgoing-boundary descriptions impose the same pole
condition in different coordinates.  In a tight-binding lead, for example,
eliminating the semi-infinite lead produces a surface Green function in
$\vsub{H}{eff}(z)$.  Choosing its retarded branch is exactly the choice of
an outgoing wave in the lead.  Conversely, substituting an outgoing ansatz at
the boundary reduces the open problem to a finite region with a
spectral-parameter-dependent boundary operator.  General versions of this
equivalence are discussed in Ref.~\cite{Hatano:2013nyz}.

This equivalence suggests a useful principle.  Non-self-adjointness is often a
representation of openness rather than a modification of the microscopic
theory.  Different representations distribute the analytic continuation
differently: in $\vsub{H}{eff}(z)$ it resides in a self-energy; in complex
scaling it resides in a deformed operator domain; in a rigged Hilbert space it
resides in the dual pairing; and in exact WKB it resides in Borel-summed
connection data.

\subsection{Resonance expansion and the role of the continuum}
\label{subsec:berggren-expansion}
% BEGIN CURATED SUBJECT INDEX
\index{operator!adjoint and self-adjoint}
\index{Berggren completeness relation}
\index{pseudostate}
% END CURATED SUBJECT INDEX

For suitable finite-range or dilation-analytic problems, one may deform the
real momentum contour into the complex plane.  The resulting Berggren-type
completeness relation has the schematic form
\begin{equation}
  \identity
  =\sum_{n\in\mathcal B}\ket{u_n}\bra{\widetilde u_n}
  +\sum_{r\in\mathcal R}\ket{u_r}\bra{\widetilde u_r}
  +\int_{L^+}\rmd k\,
  \ket{u(k)}\bra{\widetilde u(k)} .
  \label{eq:berggren-completeness}
\end{equation}
The set $\mathcal B$ contains bound states, $\mathcal R$ contains resonance
poles enclosed between the original and deformed contours, and $L^+$ is the
remaining continuum contour.  Tildes denote the dual system appropriate to
the non-self-adjoint representation.  Equation~\eqref{eq:berggren-completeness}
is not a statement that resonance vectors alone form a complete basis.  The
deformed continuum is indispensable and carries threshold and background
response
\cite{Berggren:1968zz,Michel:2009,Myo:2014ypa}.

In a finite numerical basis the continuum integral is replaced by a string of
discrete points.  The points move when the contour, scaling angle, or basis is
changed, whereas an exposed isolated pole is stable.  This contrast is the
first practical diagnostic for distinguishing a resonance from a continuum
pseudostate.
